\answer{ex:01:1}{(а)~$8\cdot10^{10}$ \nt{GLUT} --- $10$ населень Землі; широкомовних $\approx1.9\cdot10^6$ --- місто завбільшки з Відень. (б)~$\approx8\cdot10^{14}$ закінчень, із них широкомовних $\approx3\cdot10^{11}$, частка $\approx4\cdot10^{-4}$ --- у $\approx16$ разів більша за частку цих нейронів ($2.2\cdot10^{-5}$).}
\answer{ex:01:2}{(а)~$\tau_c\ln(c_0/c_*)\approx4.6\ms$. (б)~Лінія вільна при $T\ge\tau_c\ln101$: до $\approx22$ викидів за секунду замість $\approx220$.}
\answer{ex:01:3}{$V(s_k)=\gamma^{K-1-k}r$; $\delta_{\text{лампа}}=\gamma^Kr=re^{-T/\tau_\gamma}\approx0.60\,r$ за будь-якого $\Delta$, $\tau_\gamma\approx1.95\,\text{с}$.}
\answer{ex:01:4}{$n^*=93$, $47$, $25$, $12$: $n^*\approx5/\alpha$ при малих $\alpha$.}
\answer{ex:01:5}{\nt{GLUT}: $1\to10\to10^4\to10^7$ --- $\approx10^7$ нейронів, $4$ ланки, $4\ms$. \nt{DOPA}: $1\to10\to3.2\cdot10^5$ --- $\approx3\cdot10^5$ нейронів, $3$ ланки, у $30$ разів менше; той самий порядок, що й кількість \nt{DOPA}-нейронів.}
\answer{ex:01:6}{$\lambda=f_EJ_E-f_IJ_I$, звідки $J_I=37.5$; відношення струмів $f_IJ_I/f_EJ_E=0.94$; лавина ($\lambda\ge1$) при послабленні гальмування всього на $6.7\,\%$.}
\answer{ex:01:7}{$n\ge57$ пред'явлень на кожну затримку в кожному стані. Такий самий спад дала б, наприклад, неточність внутрішньої оцінки часу, що зростає з $T$.}
